\section{深度 RL 中的探索}

\subsection{从 Bandit 到 MDP}

在 MDP 中，探索问题更加复杂：
\begin{itemize}
    \item 动作不仅影响即时奖励，还影响未来的状态
    \item 需要探索\textbf{状态--动作对}而非仅是动作
    \item 有些状态本身就很难到达，需要长时间的定向探索
\end{itemize}

\subsection{基于计数的探索}

\begin{importantbox}{Count-Based Exploration}
基本思想：对访问次数少的状态--动作对给予额外的探索奖励：
$$r_{\text{explore}}(s, a) = \frac{\beta}{\sqrt{N(s, a)}}$$
在深度 RL 中，由于状态空间连续且高维，无法直接计数。常用的近似方法包括：
\begin{itemize}
    \item 哈希计数：将状态通过哈希函数映射到离散空间
    \item 密度模型：训练密度模型 $\hat{p}(s)$，用 $-\log \hat{p}(s)$ 作为新颖性度量
    \item RND（Random Network Distillation）：用随机网络的预测误差度量新颖性
\end{itemize}
\end{importantbox}

\subsection{基于好奇心的探索}

另一种思路：探索奖励 = 模型预测误差。如果对某个状态的预测误差大，说明这个状态是"新颖的"，值得继续探索。

\begin{warningbox}{好奇心陷阱}
基于预测误差的探索方法有一个著名的失败模式：\textbf{噪声电视问题}（noisy TV problem）。如果环境中存在本质上不可预测的随机性（如电视屏幕上的静态噪声），模型的预测误差永远不会降低，导致智能体被这种无意义的随机性"吸引"，无法继续有意义的探索。
\end{warningbox}

\subsection{为什么机器人和 LLM 很少从零探索}

课程在这部分给出的结论相当务实：对于机器人控制和 LLM 这类巨大 MDP，从零开始做有效探索通常在计算上不可承受。工业界和前沿研究更常见的策略是：
\begin{itemize}
    \item 使用 demonstrations 或预训练 base model 缩小搜索空间；
    \item 尽可能提供 shaped rewards；
    \item 把真正困难的探索留给更小的、结构化得更好的子问题。
\end{itemize}

\begin{figure}[H]
\centering
\includegraphics[width=0.9\textwidth]{slides-images/slide-13.jpg}
\caption{在机器人和 LLM 中，通常依赖预训练、示范和 shaped reward，而非从零探索}
\end{figure}

\subsection{本章小结}

深度 RL 中的探索需要在连续高维空间中处理新颖性度量，count-based 方法和好奇心方法各有优势和局限。

